\documentclass[12pt]{article}
\usepackage[T1]{fontenc}
\usepackage[utf8]{inputenc}
\usepackage{lmodern}
\usepackage{textcomp}
\usepackage{enumitem}
\usepackage{geometry}
\geometry{margin=1in}
\begin{document}
\begin{center}
Answer Key - Astronomy - Graduate Level Assessment 8 \
\small (Answers are ordered exactly as in the question paper.)
\end{center}

\section*{Multiple Choice}
\begin{enumerate}[label=\arabic*.]
\item \textbf{Answer:} C
\\ \textit{Options:} A) ancient heavily cratered highlands; B) dark lavas inside volcanic calderas; C) dark lavas filling older impact basins; D) the bright regions on the Moon
\\ \textit{Note:} The lunar maria are characterized as dark, basaltic plains formed by ancient volcanic activity that filled large impact basins on the Moon's surface, distinguishing them from the lighter, heavily cratered highlands.
\item \textbf{Answer:} C
\\ \textit{Options:} A) Solar eclipses would be much rarer.; B) Solar eclipses would last much longer.; C) Solar eclipses would be much more frequent.; D) Solar eclipses would not last as long.
\\ \textit{Note:} If the Moon's orbital plane were aligned with the ecliptic, the conditions for solar eclipses would be met more often during new moons, resulting in a significantly increased frequency of these events.
\item \textbf{Answer:} D
\\ \textit{Options:} A) From the layering of materials within the Earth.; B) From fossils of ancient life.; C) From the cratering history of Earth's surface.; D) From radioactive dating of rocks and meteorites.
\\ \textit{Note:} The Earth's age is determined through radioactive dating, which measures the decay of isotopes in rocks and meteorites, providing a reliable timeline that indicates the planet formed approximately 4.5 billion years ago.
\item \textbf{Answer:} C
\\ \textit{Options:} A) Centaurus; B) Cygnus; C) Cassiopeia; D) Cepheus
\\ \textit{Note:} Cassiopeia is recognized as a prominent W-shaped constellation in the northern sky, characterized by its distinctive arrangement of five bright stars that form the shape of a 'W'.
\item \textbf{Answer:} B
\\ \textit{Options:} A) Because it's climbing up a big hill.; B) Because it's in the southern hemisphere where it is winter now.; C) Because it's in the northern hemisphere where it is winter now.; D) Because one of its wheels broke.
\\ \textit{Note:} The Mars Exploration Rover Spirit is tilted towards the north because, during winter in the southern hemisphere of Mars, it positions itself to receive more sunlight by tilting its solar panels towards the sun, which is located in the northern sky.
\end{enumerate}

\section*{True / False}
\begin{enumerate}[label=\arabic*.]
\setcounter{enumi}{5}
\item \textbf{Answer:} True
\\ \textit{Options:} A) True; B) False
\\ \textit{Note:} The statement is true because the Moon is tidally locked to the Earth, meaning its rotational period on its axis matches its orbital period around the Earth, causing the same hemisphere to always face our planet.
\item \textbf{Answer:} True
\\ \textit{Options:} A) True; B) False
\\ \textit{Note:} Meteorites with high metal content are primarily composed of iron and nickel, which are characteristic of the core materials of differentiated asteroids that have undergone significant internal melting and differentiation, and their fragmentation due to collisions can lead to the ejection of these metal-rich fragments into space.
\item \textbf{Answer:} False
\\ \textit{Options:} A) True; B) False
\\ \textit{Note:} The statement is false because a neutral atom can have varying numbers of neutrons and protons, as the stability of an atom is determined by the balance of protons and neutrons, not their equality.
\item \textbf{Answer:} False
\\ \textit{Options:} A) True; B) False
\\ \textit{Note:} The statement is false because the telescope was invented in the early 1600 s, with Galileo making significant improvements to it around 1609, which laid the groundwork for modern astronomical study, not in 1709.
\item \textbf{Answer:} False
\\ \textit{Options:} A) True; B) False
\\ \textit{Note:} Meteorites do not always exhibit a fusion crust because some may break apart or disintegrate before reaching the Earth's surface, while others may not experience sufficient atmospheric heating to form a crust.
\end{enumerate}

\section*{Long Form Response}
\begin{enumerate}[label=\arabic*.]
\setcounter{enumi}{10}
\item \textbf{Answer:} Barnard's Star is significant due to its status as the second closest known star system to Earth, located approximately 5.96 light-years away. Its relative proximity allows for detailed study of stellar characteristics and motion, which has implications for understanding stellar evolution and the dynamics of our solar neighborhood. Compared to Proxima Centauri, which is slightly closer at 4.24 light-years and part of the Alpha Centauri system, Barnard's Star is a red dwarf with a much lower luminosity and mass than its counterparts, which impacts its habitability potential. Furthermore, while Proxima Centauri hosts an exoplanet in its habitable zone, Barnard's Star, despite its intriguing features, has not been confirmed to have any planets. Thus, while Barnard's Star serves as a critical point of reference in stellar studies, it highlights the diversity and complexity of our nearest stellar neighbors.
\\ \textit{Note:} correct because it accurately identifies Barnard's Star as the second closest star to Earth, emphasizes its importance for studying stellar characteristics and motion, and compares its classification as a red dwarf to Proxima Centauri, highlighting key differences in proximity and stellar type within the context of the Alpha Centauri system.
\item \textbf{Answer:} Ptolemy's use of 'circles upon circles' in his geocentric model was primarily motivated by the need to account for the observed retrograde motion of planets, where they appear to move westward against the backdrop of the stars. He introduced epicycles-smaller circles whose centers move along larger circles (deferents)-to explain this phenomenon, as the combination of these circular motions allowed for a more accurate representation of the complex paths that planets take in the sky. This approach adhered to the philosophical and observational paradigm of the time, emphasizing uniform circular motion as the most natural form of celestial movement. By employing this intricate system, Ptolemy was able to provide a mathematical framework that matched the observational data available, thereby maintaining the geocentric perspective while addressing the challenges posed by retrograde motion. Ultimately, this model exemplified the blending of empirical observations with theoretical constructs, a hallmark of ancient astronomical thought.
\\ \textit{Note:} correct because it accurately identifies Ptolemy's implementation of epicycles within his geocentric model as a solution to explaining retrograde motion, thereby demonstrating how the complex circular motions of planets could be reconciled with observational data.
\end{enumerate}

\vfill
\noindent\textit{End of Answer Key}
\end{document}